\chapter{Clustering, Regimes and Anomaly Detection}\label{ch:ml:clustering-regimes-and-anomaly-detection}

A risk dashboard shows the market in its calm regime on the Friday before a crash. On Tuesday, refitted, the same model shows
the crash beginning on Friday. Both were right about what they could know: the first was a filtered probability, computed
from the data up to Friday; the second a smoothed one, computed with Monday in hand. On the chapter's planted regimes the
filtered probability of stress on the day before a switch averages 0.03, no higher than on any calm day; the smoothed one
averages 0.35. This chapter is about learning without labels: grouping assets and days (\omterm{def:ml:clustering-regimes-and-anomaly-detection:cluster}{clustering}), describing the market
by hidden states (regime models), and finding what does not fit (\omterm{def:ml:clustering-regimes-and-anomaly-detection:anomaly}{anomaly detection}). Each is useful, and each invites a
particular mistake: clusters that are noise, regimes known only in hindsight, anomalies that are merely rare.

\section{Clustering assets and days}

\begin{definition}[Clustering, k-means, hierarchical clustering]\label{def:ml:clustering-regimes-and-anomaly-detection:cluster}
\emph{Clustering}\index{clustering} partitions objects into groups whose members are more similar to each other than to
others, without labels. \emph{k-means}\index{k-means} chooses $k$ centres and assigns each object to the nearest, alternating
the two steps to minimise the within-group sum of squared distances (Lloyd's algorithm). \emph{Hierarchical
clustering}\index{hierarchical clustering} merges the closest groups step by step into a tree and cuts it at the desired number
of groups; for assets the usual distance is $\sqrt{2(1-\rho_{ij})}$ between return series (Mantegna, 1999), the one behind
hierarchical risk parity (Book~7, chapter~26).
\end{definition}

The chapter clusters the names of the firm's synthetic market (Book~7's \texttt{firm.synthmkt}) that stay listed for five
years, whose returns load on a market factor, ten planted industry factors, style factors and fat-tailed specific noise. The
market factor is removed first (each name's return minus its beta times the equal-weighted average), since it would dominate
every correlation. \Cref{fig:ml:reg:clusters} scores \omterm{def:ml:clustering-regimes-and-anomaly-detection:cluster}{k-means} and average-linkage \omterm{def:ml:clustering-regimes-and-anomaly-detection:cluster}{hierarchical clustering} by the adjusted Rand
index against the planted industries (1 is perfect, 0 is chance), for windows of three months to five years and for two
strengths of the industry factor: the default 8\% a year, and 15\%.

\begin{omfigure}
\begin{tikzpicture}
\begin{semilogxaxis}[width=11cm,height=6cm,xlabel={window (trading days)},ylabel={adjusted Rand index},
  tick label style={font=\small},label style={font=\small},xmin=50,xmax=1500,ymin=-0.05,ymax=1.05,
  xtick={63,252,1260},xticklabels={63,252,1\,260},grid=major,grid style={gray!25},legend columns=2,legend cell align=left,
  legend style={font=\footnotesize,at={(0.5,-0.25)},anchor=north,draw=none,fill=none,
  /tikz/every even column/.append style={column sep=6pt}},table/col sep=comma]
\addplot[omThm,very thick,mark=*] table[x=window,y=km15]{figdata/ml/20-clustering-regimes-and-anomaly-detection/clusters.csv};
\addplot[omDef,very thick,mark=square*] table[x=window,y=hier15]{figdata/ml/20-clustering-regimes-and-anomaly-detection/clusters.csv};
\addplot[omThm!50,thick,dashed,mark=o,mark options={solid}] table[x=window,y=km08]{figdata/ml/20-clustering-regimes-and-anomaly-detection/clusters.csv};
\addplot[omDef!50,thick,dashed,mark=square,mark options={solid}] table[x=window,y=hier08]{figdata/ml/20-clustering-regimes-and-anomaly-detection/clusters.csv};
\legend{k-means (industry 15\%),hierarchical (industry 15\%),k-means (industry 8\%),hierarchical (industry 8\%)}
\end{semilogxaxis}
\end{tikzpicture}

\omcaption{Agreement of return-based clusters with the planted industries, by the length of the return history and the
strength of the industry factor. Data: \texttt{ml\_regimes.clustering}.}\label{fig:ml:reg:clusters}
\end{omfigure}

With strong industries, \omterm{def:ml:clustering-regimes-and-anomaly-detection:cluster}{clustering} recovers them from five years of data (index 0.99 and 1.00) and half of them from one
year; with the default strength, which is closer to what specific risk leaves of real industry effects, it finds almost
nothing in a year (0.06) and 0.15 to 0.21 in five. The half-sample stability score measures the same thing without labels: the
index between \omterm{def:ml:clustering-regimes-and-anomaly-detection:cluster}{clusterings} of the two halves of a year is 0.32 to 0.35 with strong industries and 0.06 to 0.08 with weak ones.
A \omterm{def:ml:clustering-regimes-and-anomaly-detection:cluster}{clustering} that does not survive a split of its own sample is describing noise, and a desk that uses it to group risk or
build hierarchical risk parity is rebalancing on noise.

\section{Regime models}

\begin{definition}[Hidden Markov model, regime-switching model, Viterbi algorithm]\label{def:ml:clustering-regimes-and-anomaly-detection:hmm}
A \emph{hidden Markov model}\index{hidden Markov model} (HMM) has an unobserved state that follows a Markov chain (Book~4,
chapter~8) and observations whose distribution depends on the current state; it is fitted by expectation-maximisation (Book~4,
chapter~19), the Baum--Welch algorithm (Rabiner, 1989). A \emph{regime-switching model}\index{regime-switching model} is an HMM
for economic series, each state a regime with its own mean, volatility or dynamics (Hamilton, 1989; Ang and Timmermann, 2012).
The \emph{Viterbi algorithm}\index{Viterbi algorithm} finds the single most likely path of states given all the observations,
by dynamic programming.
\end{definition}

The chapter's index switches between a calm regime (daily mean 5 basis points, volatility 0.8\%) and a stress regime (mean
$-10$ basis points, volatility 2.5\%), staying in each with probability 0.99 and 0.95 a day; twenty years are simulated, and
stress takes a fifth of the time. A two-state Gaussian HMM fitted on the first ten years finds calm with volatility 0.80\% and
stress with 2.60\%, staying probabilities 0.988 and 0.958. A three-state model adds 5.7 to the log-likelihood with six more
parameters, less than the Bayesian information criterion's penalty of 23.5, and spends its third state splitting stress in two
(2.28\% and 2.65\%).

\section{Filtering against smoothing: what a regime model knows in real time}

An HMM gives three answers about the state at date $t$. The \emph{filtered} probability uses the data up to $t$ (the forward
recursion of \cref{lst:ml:reg:hmm}); the \emph{smoothed} probability uses all the data, including what came after $t$ (the
backward pass); the Viterbi path is the best sequence given all the data. Only the first existed at $t$. On the ten test years,
the filtered estimate (stress when its probability exceeds one half) is right on 95.9\% of days, the smoothed on 97.7\% and the
Viterbi path on 97.5\%; a rule on the trailing 20-day volatility is right on 89.8\%. After each of the 20 switches into stress,
the filter says stress after a median of one day (mean 1.5, one short spell missed), the smoothed estimate on the day itself
(mean 0.8). \Cref{fig:ml:reg:switch} shows one switch: the smoothed probability rises before the regime changes, because the
days after the switch make the days before it look like stress.

\begin{omfigure}
\begin{tikzpicture}
\begin{axis}[width=12cm,height=6cm,xlabel={days from a switch into stress},ylabel={probability of stress},
  tick label style={font=\small},label style={font=\small},xmin=-30,xmax=49,ymin=-0.02,ymax=1.05,grid=major,
  grid style={gray!25},legend columns=3,legend cell align=left,legend style={font=\footnotesize,at={(0.5,-0.25)},anchor=north,
  draw=none,fill=none,/tikz/every even column/.append style={column sep=6pt}},table/col sep=comma]
\addplot[gray!60,fill=gray!15,const plot,draw=none] table[x=day,y=state]{figdata/ml/20-clustering-regimes-and-anomaly-detection/switch.csv} \closedcycle;
\addplot[omThm,very thick] table[x=day,y=filtered]{figdata/ml/20-clustering-regimes-and-anomaly-detection/switch.csv};
\addplot[omDef,very thick,dashed] table[x=day,y=smoothed]{figdata/ml/20-clustering-regimes-and-anomaly-detection/switch.csv};
\legend{true stress regime,filtered (known that day),smoothed (known at the end)}
\end{axis}
\end{tikzpicture}

\omcaption{Filtered and smoothed probabilities of stress around one switch in the test years (the first switch the filter took
at least two days to see). Data: \texttt{ml\_regimes} (\texttt{fig\_regimes.py}).}\label{fig:ml:reg:switch}
\end{omfigure}

The difference is a look-ahead bias (Book~7, chapter~3) that no timestamp reveals. On the day before a switch, the filtered
probability of stress averages 0.032, below its average of 0.041 over all calm days: the filter cannot see a switch coming. The
smoothed probability averages 0.353 on those days against 0.019 on calm days. A backtest that scales exposure to a target
volatility using the regime probabilities earns a Sharpe ratio of 1.00 with yesterday's filtered probability carried forward
by the transition matrix, and 1.21 with the same day's smoothed probability; buying and holding earns 0.91. A regime model
used for trading or risk must be run as a filter, refitted only on data available at each date, and its statistics
(the volatility of a regime, its average duration) computed from filtered probabilities.

\section{Anomaly detection}

\begin{definition}[Anomaly detection, isolation forest, precision--recall curve]\label{def:ml:clustering-regimes-and-anomaly-detection:anomaly}
\emph{Anomaly detection}\index{anomaly detection} scores observations by how unlike the bulk of the data they are, without
labels for the anomalies. An \emph{isolation forest}\index{isolation forest} builds random trees that split on random features
at random values; points that are isolated after few splits are anomalous (Liu, Ting and Zhou, 2008). The
\emph{precision--recall curve}\index{precision--recall curve} plots, as the alarm threshold moves, the share of alarms that are
true (precision) against the share of true cases that raise an alarm (recall); when positives are rare it is more informative
than the ROC curve (Davis and Goadrich, 2006).
\end{definition}

Book~9 (chapter~29) built surveillance detectors for spoofing on account-days: 100 planted episodes among 11\,500 legitimate
account-days, including market makers who cancel almost everything and deep-book liquidity providers whose large orders rarely
fill. \Cref{tab:ml:reg:anomaly} adds two unsupervised scores on five features of the same account-days (the log counts of small
and large orders, their fill rates, and the share of large cancellations that follow a fill on the other side): an \omterm{def:ml:clustering-regimes-and-anomaly-detection:anomaly}{isolation
forest} (\cref{lst:ml:reg:iforest}) and the reconstruction error of a small \omterm{def:ml:cross-sectional-deep-models:ae}{autoencoder} (chapter~9).

\begin{table}[htbp]
\centering\small
\begin{tabular}{@{}lcc@{}}
\toprule
detector & recall at 1\% false positives & precision\\
\midrule
order-to-trade ratio (rule) & 0.00 & 0.00\\
fill-rate gap (rule) & 0.37 & 0.24\\
cancellations after a fill (rule) & 0.51 & 0.31\\
gap $\times$ cancellations (rule) & 0.77 & 0.40\\
\omterm{def:ml:clustering-regimes-and-anomaly-detection:anomaly}{isolation forest} & 1.00 & 0.47\\
\omterm{def:ml:cross-sectional-deep-models:ae}{autoencoder} & 0.32 & 0.22\\
\bottomrule
\end{tabular}
\caption{Spoofing detectors on 11\,600 account-days (100 planted episodes), at the threshold that flags 1\% of legitimate
account-days. Data: \texttt{ml\_regimes.anomalies}.}\label{tab:ml:reg:anomaly}
\end{table}

\section{Surveillance detectors}

The \omterm{def:ml:clustering-regimes-and-anomaly-detection:anomaly}{isolation forest} finds every planted episode at the 1\% false-positive rate, and its precision, 0.47, is the most any
detector can have there: 115 legitimate account-days are flagged whatever the detector, against at most 100 true ones. The
\omterm{def:ml:cross-sectional-deep-models:ae}{autoencoder} does worse than the rules. Three cautions keep this result in its place. The forest was given the features the rules'
authors chose, which encode what spoofing looks like; on raw order logs it would find the unusual, not the manipulative. The
planted spoofers are a small, distinct population; when a manipulative pattern is common, or when legitimate traders who look
like it are many, an unsupervised score ranks them together. And precision at a fixed false-positive rate is what a compliance
team lives with: a detector with perfect recall still hands them more innocent account-days than guilty ones. Unsupervised
scores earn their place as a second net under the rules, reviewed by people, and as a way to find patterns no rule names yet.

\begin{method}[Unsupervised learning on market data]\label{met:ml:reg:method}
\begin{enumerate}
\item Remove the common factor before \omterm{def:ml:clustering-regimes-and-anomaly-detection:cluster}{clustering} returns, and report a stability score across sub-samples with every
\omterm{def:ml:clustering-regimes-and-anomaly-detection:cluster}{clustering}.
\item Fit regime models on data available at each date and use filtered probabilities only; report smoothed results as
descriptions of history, never as backtests.
\item Choose the number of regimes by an out-of-sample likelihood or an information criterion, and check that the regimes differ
in something the desk acts on.
\item Evaluate anomaly scores like detectors: recall and precision at the false-positive rate reviewers can handle, against
the rules, on data with planted and look-alike cases.
\end{enumerate}
\end{method}

\section{Tutorial: which regime are we in?}\label{tut:ml:clustering-regimes-and-anomaly-detection:tut}

\begin{tutorial}
\textbf{Goal.} Cluster the synthetic names, fit and run a regime model as a filter and as a smoother, and score unsupervised
detectors against the surveillance rules. \textbf{End state:} \cref{fig:ml:reg:clusters,fig:ml:reg:switch},
\cref{tab:ml:reg:anomaly}.
\begin{tutsteps}
\item \textbf{\omterm{def:ml:clustering-regimes-and-anomaly-detection:cluster}{Hierarchical clustering} on the correlation distance.}
\omcode{code/firm/regimes/firm_regimes.py}{43}{50}{Average linkage on Mantegna's distance.}\label{lst:ml:reg:hier}
\item \textbf{Filtering and smoothing.}
\omcode{code/firm/regimes/firm_regimes.py}{84}{109}{The forward recursion (filter) and the backward pass (smoother).}\label{lst:ml:reg:hmm}
\item \textbf{An \omterm{def:ml:clustering-regimes-and-anomaly-detection:anomaly}{isolation forest}.}
\omcode{code/firm/regimes/firm_regimes.py}{161}{166}{Isolation-forest anomaly scores.}\label{lst:ml:reg:iforest}
\item \textbf{Run} \texttt{ml\_regimes.clustering()}, \texttt{regimes()}, \texttt{vol\_targeting()}, \texttt{anomalies()} and
\texttt{fig\_regimes.py}.
\end{tutsteps}
\textbf{What to change next.} Refit the HMM every year on an expanding window and filter the next year; double the number of
legitimate deep-book providers and rerun the anomaly scores.
\end{tutorial}

\section{Build: clustering, regimes and anomaly scores}\label{bld:ml:clustering-regimes-and-anomaly-detection:regimes}

\begin{build}
\bfield{Purpose} Unsupervised tools whose outputs are known in real time and scored like detectors.
\bfield{Interface} \texttt{market\_residuals}, \texttt{corr\_distance}, \texttt{kmeans\_clusters}, \texttt{hier\_clusters},
\texttt{ari}, \texttt{stability}; \texttt{regime\_series}; \texttt{GaussianHMM(k)} with \texttt{fit}, \texttt{filter},
\texttt{smooth}, \texttt{viterbi}; \texttt{iforest\_scores}, \texttt{autoencoder\_scores}; with Book~9's
\texttt{firm.surveil.tpr\_at\_fpr}.
\bfield{Rules} Filtered probabilities only in anything that trades or reports risk; every \omterm{def:ml:clustering-regimes-and-anomaly-detection:cluster}{clustering} with its stability score;
anomaly scores with recall and precision at a fixed false-positive rate.
\bfield{Acceptance tests} \texttt{code/firm/regimes/tests/}: \omterm{def:ml:clustering-regimes-and-anomaly-detection:cluster}{clustering} recovers well-separated planted groups; the HMM recovers
planted parameters and its filter uses no future data (it is unchanged when later data change); smoothing and filtering agree
on the last day; Viterbi matches brute force on a short series; the \omterm{def:ml:clustering-regimes-and-anomaly-detection:anomaly}{isolation forest} ranks planted outliers first.
\bfield{Stretch} Regimes with their own dynamics (autoregressive or GARCH within states); online refitting; \omterm{def:ml:clustering-regimes-and-anomaly-detection:cluster}{clustering} days
rather than assets.
\end{build}

\begin{omsources}
\item J.~D.~Hamilton, ``A new approach to the economic analysis of nonstationary time series and the business cycle'',
\emph{Econometrica} 57(2), 1989.
\item A.~Ang and A.~Timmermann, ``Regime changes and financial markets'', \emph{Annual Review of Financial Economics} 4, 2012.
\item L.~R.~Rabiner, ``A tutorial on hidden Markov models and selected applications in speech recognition'',
\emph{Proceedings of the IEEE} 77(2), 1989.
\item F.~T.~Liu, K.~M.~Ting and Z.-H.~Zhou, ``Isolation forest'', \emph{IEEE International Conference on Data Mining}, 2008.
\item S.~Lloyd, ``Least squares quantization in PCM'', \emph{IEEE Transactions on Information Theory} 28(2), 1982.
\item R.~N.~Mantegna, ``Hierarchical structure in financial markets'', \emph{European Physical Journal B} 11, 1999.
\item J.~Davis and M.~Goadrich, ``The relationship between precision-recall and ROC curves'', \emph{ICML}, 2006.
\end{omsources}

\section{Exercises}

\begin{exercise}[$\star$]\label{exo:ml:clustering-regimes-and-anomaly-detection:1}
Two names have return correlation 0.5. What is their Mantegna distance? What distance do perfectly correlated and uncorrelated
names have?
\end{exercise}

\begin{exercise}[$\star$]\label{exo:ml:clustering-regimes-and-anomaly-detection:2}
With staying probabilities 0.988 and 0.958, what are the expected durations of the calm and stress regimes, and the long-run
share of time in stress?
\end{exercise}

\begin{exercise}[$\star$]\label{exo:ml:clustering-regimes-and-anomaly-detection:3}
At a 1\% false-positive rate on 11\,500 legitimate account-days, what is the highest precision any detector can reach with 100
true cases?
\end{exercise}

\begin{exercise}[$\star\star$]\label{exo:ml:clustering-regimes-and-anomaly-detection:4}
Why is the filtered probability of stress on the day before a switch no higher than on an average calm day?
\end{exercise}

\begin{exercise}[$\star\star$]\label{exo:ml:clustering-regimes-and-anomaly-detection:5}
Why must the market factor be removed before \omterm{def:ml:clustering-regimes-and-anomaly-detection:cluster}{clustering} returns?
\end{exercise}

\begin{exercise}[$\star\star$]\label{exo:ml:clustering-regimes-and-anomaly-detection:6}
\emph{Find the flaw.} ``Our two-regime model's stress probability, fitted on 2000--2024, turned up weeks before each of the
three big drawdowns, so it is an early-warning signal.''
\end{exercise}

\begin{exercise}[$\star\star\star$]\label{exo:ml:clustering-regimes-and-anomaly-detection:7}
\emph{Coding.} Compare the three-state model with the two-state one on the test years by the log-likelihood of the test data
under each fitted model (filtering forward from the fitted parameters). Which wins out of sample?
\end{exercise}

\begin{exercise}[$\star\star\star$]\label{exo:ml:clustering-regimes-and-anomaly-detection:8}
Derive the forward recursion $\alpha_t(j) = \big(\sum_i\alpha_{t-1}(i)P_{ij}\big)b_j(x_t)$ for the joint probability of the
observations up to $t$ and the state at $t$, and show that normalising it gives the filtered probability.
\end{exercise}

\section{Problem: Which Regime Are We In?}

\begin{problem}[{Weekend problem --- filtered, smoothed, and rare}]\label{pb:ml:clustering-regimes-and-anomaly-detection:1}
The chapter's clusters, regimes and detectors.

\medskip\noindent\textbf{Part I --- Clusters.}
\begin{enumerate}
\item How well do clusters recover industries, by window and strength?
\item What does the stability score say, and why is it the one to report on real data?
\item What would you use clusters for on a desk, and what not?
\item Which \omterm{def:ml:clustering-regimes-and-anomaly-detection:cluster}{clustering} method would you choose, and why?
\end{enumerate}

\medskip\noindent\textbf{Part II --- Regimes.}
\begin{enumerate}[resume]
\item What does the fitted two-state model find, against the planted regimes?
\item Why not three states?
\item What are the accuracies of the four regime estimates on the test years?
\item What are the delays after switches into stress?
\end{enumerate}

\medskip\noindent\textbf{Part III --- Hindsight.}
\begin{enumerate}[resume]
\item What do the filtered and smoothed probabilities say on the day before a switch?
\item How much does the smoothed probability flatter a volatility-targeting backtest?
\item How should regime statistics be computed?
\item Where else does smoothing hide in a research pipeline?
\end{enumerate}

\medskip\noindent\textbf{Part IV --- The verdict.}
\begin{enumerate}[resume]
\item State the \emph{named result}: the filtered regime's accuracy and detection delay against the smoothed one, and the anomaly
detectors' precision at a fixed false-positive rate against the rule-based detectors.
\item What is a regime model for, if not prediction?
\item Would you let the \omterm{def:ml:clustering-regimes-and-anomaly-detection:anomaly}{isolation forest} replace the rules?
\item How would you review the \omterm{def:ml:clustering-regimes-and-anomaly-detection:anomaly}{isolation forest}'s alarms?
\item What would you monitor to know that a regime model has broken?
\item How do clusters feed hierarchical risk parity, and what happens when they are noise?
\item What labels would turn the anomaly problem into a supervised one, and at what cost?
\item In one sentence: what does a model know about today?
\end{enumerate}
\end{problem}

\section{Interview questions}

\begin{interviewq}[$\star$ \iqroles{researcher}]\label{iq:ml:clustering-regimes-and-anomaly-detection:1}
What is the difference between filtered and smoothed state probabilities in an HMM?
\end{interviewq}

\begin{interviewq}[$\star\star$ \iqroles{researcher, mle}]\label{iq:ml:clustering-regimes-and-anomaly-detection:2}
How would you cluster 500 stocks by their returns, and how would you know the clusters are real?
\end{interviewq}

\begin{interviewq}[$\star\star$ \iqroles{researcher}]\label{iq:ml:clustering-regimes-and-anomaly-detection:3}
Explain the Baum--Welch algorithm in a few sentences.
\end{interviewq}

\begin{interviewq}[$\star\star$ \iqroles{mle}]\label{iq:ml:clustering-regimes-and-anomaly-detection:4}
How does an \omterm{def:ml:clustering-regimes-and-anomaly-detection:anomaly}{isolation forest} work, and what are its blind spots?
\end{interviewq}

\begin{interviewq}[$\star\star$ \iqroles{researcher, trader}]\label{iq:ml:clustering-regimes-and-anomaly-detection:5}
A regime-switching overlay improved a strategy's backtested Sharpe ratio from 0.9 to 1.2. What do you check?
\end{interviewq}

\begin{interviewq}[$\star\star\star$ \iqroles{mle}]\label{iq:ml:clustering-regimes-and-anomaly-detection:6}
Design an anomaly-detection system for trade surveillance. How do you evaluate it without many labels?
\end{interviewq}
